
\chapter{Ao lado: o dispositivo}

\noindent\emph{Verbete acrescentado em 30 de outubro de 2026, em diálogo com o operador, depois de ele ratificar a
primeira leitura do capítulo anterior. Nada do que está escrito acima muda; isto fica ao lado.}

A primeira leitura --- ``a gramática é a do registro público'' --- foi ratificada: \emph{``sim, está certa''}. E
você pediu para definir a palavra que ela usa:
\citacao{MAs [sic] vamos tentar definir ``dispositivo'', porque no meu mundo jurídico trata-se do sujeito da
segurança jurídica e em linguística é a ontologia do nome o que eu defino como o operador de memória, faz sentido
pra vc assim?}{--- o operador, 30/10/2026; verbatim por sha256 \texttt{cb7a3e83e580a548}}

Perguntado se ``sujeito'' queria dizer o portador ou o produtor, você respondeu:
\citacao{ser o sujeito da segurança jurídica signiifca [sic] ser o objeto que ela protege ou que a operação criva,
afixa como estado de representação}{--- o operador, 30/10/2026; verbatim por sha256 \texttt{fb8a9adb5991e589}}

Os dois sentidos são um só, em dois momentos: primeiro a operação criva e afixa; depois a segurança jurídica
protege o que ficou afixado. E ``criva'' é a palavra do kernel: uma peneira é uma projeção --- muitos entram, um
estado sai, o que não passa não volta. É a diferença, já marcada neste memorial, entre o separante de Tomita, que
não perde nada, e o separador de Landauer, que perde e cobra. O dispositivo é o que passa pela peneira; a peneira
cobra; e o capítulo 8 disse quanto: a entropia de decidir em que metade caiu, a Meia-Nat.

\eqbox{\textbf{Dispositivo.} O objeto que a operação criva e afixa como estado de representação, e que a segurança
jurídica, por isso, protege. No direito, a única parte da decisão que faz coisa julgada: reler não a altera. Na
linguística, a ontologia do nome: o nome é o que permanece idêntico sob releitura. No kernel, o ponto fixo da
operação idempotente, $E^2 = E$. \textbf{O operador de memória é a operação que criva}: devolve o dispositivo e
recusa o resto. O que ela recusa é a fundamentação, que pode variar sem mudar o estado afixado.}

\noindent Ratificado pelo operador em 30/10/2026: \emph{``ratificado, é isso, e é lindo''} (verbatim por sha256
\texttt{HASH_RATIFICACAO}). No kernel, os dois mundos já se tocavam antes do verbete: a estabilidade do colapso é
provada por \texttt{res\_judicata\_is\_terminal}, e a leitura do observador por \texttt{observing\_adds\_nothing}.
A coisa julgada e a idempotência são o mesmo teorema, lido em duas línguas. \texttt{[REAL --- os termos estão no
kernel; a identificação é cunhagem do operador, ONTO]}

